\documentclass[11pt]{article}
\usepackage[margin=1in]{geometry}
\usepackage{amsmath,amssymb,amsthm,mathtools}
\usepackage{graphicx}
\usepackage[hidelinks]{hyperref}
\usepackage{microtype}

\newtheorem{theorem}{Theorem}
\newtheorem{proposition}{Proposition}
\newtheorem{corollary}{Corollary}
\newcommand{\R}{\mathbb R}
\newcommand{\C}{\mathbb C}
\newcommand{\one}{\mathbf 1}
\newcommand{\op}{\mathrm{op}}
\newcommand{\Wtwo}{\mathcal W_2}

\title{Constant Spectral Data Do Not Control\\
Matrix Schr\"odinger Maximal Regularity}
\author{Run \texttt{fa\_banach\_001} -- lane 1}
\date{27 August 2026}

\begin{document}
\maketitle

\noindent\textbf{Status.} Candidate full counterexample, likely valid, to
pointwise-spectral reverse-H\"older criteria for the system maximal estimate.
The packet also gives a sign-free positive repair under comparable
eigenvalues. The source's broader search for a geometry-sensitive matrix
condition remains open.

\section{Source question}

Bailey studies the class $\Wtwo$ associated with
$H=-\Delta I_m+|V|$ and asks for a reverse-H\"older-type condition implying
boundedness of $|V|(|V|-\Delta)^{-1}$ when $m>1$
\cite[Section~5.3, printed p.~42]{Bailey}. The exact source passage is shown
below.

\begin{center}
\includegraphics[width=0.96\textwidth,trim=0 104 0 8,clip]{figures/open_problem_crop.png}
\end{center}

For a measurable Hermitian positive-semidefinite matrix potential
$W:\R^n\to\mathcal L(\C^m)$, put
\[
 H=-\Delta I_m+W.
\]
The assertion $W\in\Wtwo$ means
\begin{equation}\label{eq:w2}
 \|Wu\|_2+\|\Delta u\|_2\le C\|Hu\|_2,
 \qquad u\in D(H).
\end{equation}

A first interpretation of the reverse-H\"older-type condition is to apply
the scalar condition to $\|W(x)\|_{\op}$, a Schatten norm, the trace, or the
ordered eigenvalue functions. The counterexample below rules out every
condition depending only on this pointwise spectral data.

\section{Proof intuition}

The scalar obstruction is the unit-ball potential
$v=\one_{B(0,1)}$ in $\R^3$. It has a bounded positive zero-energy solution
for $-\Delta+v$. Cutting this solution off at radius $R$ leaves a fixed
potential term but creates only an $O(R^{-1/2})$ residual.

To make the matrix example spectrally featureless, place $v$ in one diagonal
channel and its complement $1-v$ in the other. The resulting matrix is a
rank-one orthogonal projection at every point. Its eigenvalues, trace,
operator norm, and all Schatten norms are constant, but the bad scalar
channel survives unchanged. Thus the obstruction is the spatial motion of
the eigenspaces, information that pointwise eigenvalues cannot see.

\section{Full constant-spectrum counterexample}

\begin{theorem}\label{thm:negative}
There is a bounded measurable real-symmetric potential
$W:\R^3\to\mathcal L(\R^2)$ such that, almost everywhere,
\[
 W(x)^2=W(x),\qquad \operatorname{rank}W(x)=1,\qquad
 \operatorname{spec}W(x)=\{0,1\},
\]
but $W\notin\Wtwo$.

Consequently, both ordered eigenvalues are constant, every Schatten norm of
$W(x)$ is identically one, and
\[
 \operatorname{tr}(W(x)^k)=1\qquad(k=1,2,\ldots).
\]
In particular, no hypothesis depending only on pointwise spectral data and
satisfied by the constant spectrum $\{0,1\}$ can imply the matrix
$\Wtwo$ estimate.
\end{theorem}

\begin{proof}
Let $B=B(0,1)$, $v=\one_B$, and set
\begin{equation}\label{eq:projection}
 W(x)=
 \begin{pmatrix}
  1-v(x)&0\\
  0&v(x)
 \end{pmatrix}.
\end{equation}
Since $v^2=v$, the matrix in \eqref{eq:projection} is a rank-one orthogonal
projection at every point away from the irrelevant boundary sphere.
All asserted spectral identities follow immediately.

It remains to disprove \eqref{eq:w2}. Put $a=1-\tanh1$ and define
\[
 h(r)=
 \begin{cases}
  \dfrac{\sinh r}{r\cosh1},&0<r\le1,\\[6pt]
  1-\dfrac{a}{r},&r\ge1,
 \end{cases}
 \qquad h(0)=\dfrac1{\cosh1}.
\]
At $r=1$, both formulas have value $\tanh1$ and derivative
$a=(\cosh1-\sinh1)/\cosh1$. Hence no surface delta term occurs, and direct
radial differentiation gives
\begin{equation}\label{eq:zero}
 (-\Delta+v)h=0
\end{equation}
distributionally on $\R^3$.

Choose a fixed smooth radial cutoff $\eta$ equal to one on $[0,1]$ and zero
on $[2,\infty)$, and put
\[
 \eta_R(x)=\eta(|x|/R),\qquad
 u_R=\eta_Rh,\qquad U_R=(0,u_R),\qquad R>2.
\]
The matching just checked implies $u_R\in H^2(\R^3)$, hence
$U_R\in D(H)$. Since the second channel of $W$ is $v$,
\begin{equation}\label{eq:fixed}
 \|WU_R\|_2=\|vu_R\|_2=\|vh\|_2>0
\end{equation}
is independent of $R$. On the annulus $R<|x|<2R$, equation
\eqref{eq:zero} yields
\[
 HU_R=
 \left(0,-2\nabla\eta_R\cdot\nabla h-h\Delta\eta_R\right).
\]
There $h=O(1)$, $|\nabla h|=O(R^{-2})$,
$|\nabla\eta_R|=O(R^{-1})$, and
$|\Delta\eta_R|=O(R^{-2})$. The residual is therefore $O(R^{-2})$ on a
set of volume $O(R^3)$, so
\begin{equation}\label{eq:decay}
 \|HU_R\|_2=O(R^{-1/2})\longrightarrow0.
\end{equation}
Equations \eqref{eq:fixed}--\eqref{eq:decay} contradict
$\|WU\|_2\le C\|HU\|_2$, proving $W\notin\Wtwo$.
\end{proof}

\begin{corollary}\label{cor:spectral}
Let $\mathfrak P$ be any pointwise-unitarily-invariant condition on
positive-semidefinite matrix fields that depends only on their ordered
eigenvalue functions. If $\mathfrak P$ holds for the constant eigenvalue list
$(0,1)$, then $\mathfrak P$ does not imply membership in $\Wtwo$.
\end{corollary}

\begin{proof}
The potential \eqref{eq:projection} has that eigenvalue list everywhere and
fails $\Wtwo$ by Theorem~\ref{thm:negative}.
\end{proof}

This includes the Bochner interpretation
\[
 \left(\frac1{|Q|}\int_Q\|W(x)\|^q\,dx\right)^{1/q}
 \le \frac{C}{|Q|}\int_Q\|W(x)\|\,dx,
\]
for any unitarily invariant matrix norm: both sides are the same constant.

\section{A sign-free positive repair}

The counterexample shows that some control of eigenspace geometry or of all
quadratic-form directions is essential. One clean spectral repair is to
prevent the eigenvalues from separating.

\begin{theorem}\label{thm:positive}
Let $W:\R^n\to\mathcal L(\C^m)$ be measurable and positive semidefinite.
Write
\[
 \lambda(x)=\lambda_{\min}(W(x)),\qquad
 \Lambda(x)=\lambda_{\max}(W(x)).
\]
If, for some $q\ge2$ and $\kappa\ge1$,
\[
 \lambda\in RH_q,\qquad \Lambda\le\kappa\lambda\quad\text{a.e.},
\]
then $W\in\Wtwo$. The constant depends only on the scalar
$RH_q$ data, $n$, and $\kappa$. No off-diagonal sign condition is required,
and complex Hermitian potentials are allowed.
\end{theorem}

\begin{proof}
Let $h=-\Delta+\lambda$. Bailey's scalar theorem
\cite[Theorem~5.1]{Bailey}, quoting Auscher--Ben Ali, gives uniformly for
$\varepsilon>0$,
\begin{equation}\label{eq:scalar}
 \|\lambda(h+\varepsilon)^{-1}g\|_2\le C_s\|g\|_2.
\end{equation}
For $F\in L^2(\R^n;\C^m)$, the vector Kato inequality and
$W\ge\lambda I_m$ give the pointwise resolvent domination
\begin{equation}\label{eq:dom}
 |(H+\varepsilon)^{-1}F|
 \le (h+\varepsilon)^{-1}|F|.
\end{equation}
This uses no cooperativity or entrywise sign assumption.

Since $|Wz|\le\Lambda|z|\le\kappa\lambda|z|$, equations
\eqref{eq:scalar}--\eqref{eq:dom} imply
\[
 \|W(H+\varepsilon)^{-1}F\|_2\le\kappa C_s\|F\|_2.
\]
For $u\in D(H)$, take $F=(H+\varepsilon)u$, then let
$\varepsilon\downarrow0$. This yields
$\|Wu\|_2\le\kappa C_s\|Hu\|_2$. Finally,
$-\Delta u=Hu-Wu$, which proves \eqref{eq:w2}.
\end{proof}

\section{The geometry-sensitive frontier}

Davey and Isralowitz define a matrix reverse-H\"older class
$\mathcal B_q$ by requiring
\[
 \left(\frac1{|Q|}\int_Q\langle W(x)e,e\rangle^q\,dx\right)^{1/q}
 \le C\left\langle\left(\frac1{|Q|}\int_QW(x)\,dx\right)e,e\right\rangle
\]
uniformly over cubes $Q$ and constant vectors $e$
\cite[Definition~2.1]{DI}. This condition detects eigenspace motion and is
not contradicted by Theorem~\ref{thm:negative}: for $e=(0,1)$ the quadratic
form is $\one_B$, and for $e=(1,0)$ it is $\one_{\R^3\setminus B}$; neither
belongs to scalar $RH_q$.

Addona, Leone, Lorenzi and Rhandi prove broader $L^p$ maximal estimates under
comparable eigenvalues, a reverse-H\"older smallest eigenvalue, and
nonpositive off-diagonal entries \cite[Theorem~6.1]{ALLR}.
Theorem~\ref{thm:positive} removes the sign assumption at $p=2$.

Thus the exact remaining frontier is geometry-sensitive: for example,
whether membership in the directional matrix class $\mathcal B_2$ alone
implies $\Wtwo$ is not settled by this packet.

\section{Verification, search bounds, and recommendation}

The script \texttt{code/verify\_cutoff\_counterexample.py} checks the
projection identities, the matching conditions at $r=1$, and the exact radial
cutoff residual. For radii from $8$ to $8192$,
$\sqrt R\,\|HU_R\|_2$ stabilizes near $23.206$, while
$\|WU_R\|_2=1.4650804067$ remains fixed. The proof itself is analytic.

A bounded search through 27 August 2026 covered the run indexes, the exact
source question, arXiv and web queries for matrix-valued Schr\"odinger
maximal regularity, projection-valued potentials, constant spectrum, and
matrix reverse-H\"older classes, plus source-level inspection of
arXiv:2207.05790 and arXiv:2401.00479. No occurrence of the displayed
constant-spectrum counterexample or an equivalent $\Wtwo$ obstruction was
found. Novelty confidence is moderate because the construction is elementary.

Human review should check the interpretation boundary in
Corollary~\ref{cor:spectral}; the analytic failure of $\Wtwo$ is explicit.
The packet should be treated as a full negative result for pointwise-spectral
criteria, not as a characterization of all matrix reverse-H\"older
hypotheses.

\begin{thebibliography}{9}
\bibitem{Bailey}
J.~Bailey,
\emph{The Kato Square Root Problem for Divergence Form Operators with
Potential}, J. Fourier Anal. Appl. 26 (2020), article 46;
arXiv:1812.10196.

\bibitem{DI}
B.~Davey and J.~Isralowitz,
\emph{Exponential Decay Estimates for Fundamental Matrices of Generalized
Schr\"odinger Systems}, arXiv:2207.05790 (2022), Definition~2.1.

\bibitem{ALLR}
D.~Addona, V.~Leone, L.~Lorenzi, and A.~Rhandi,
\emph{$L^p$ Maximal Regularity for Vector-Valued Schr\"odinger Operators},
arXiv:2401.00479 (2023).
\end{thebibliography}

\end{document}
